\begingroup
\color{borrador2}
\textcolor{borrador2}{La siguiente figura propone una alternativa que separa las tres familias de generación y muestra los siete orígenes de la colección resultante.}

\begin{figure}[H]
  \centering
  \resizebox{0.96\textwidth}{!}{%
    \begin{tikzpicture}[
      nodo/.style={draw=borrador2, thick, rounded corners=3pt, fill=white,
                   text=borrador2, align=center, font=\small,
                   minimum width=3.45cm, minimum height=1.05cm},
      resultado/.style={nodo, fill=borrador2!9, font=\footnotesize\ttfamily},
      flecha/.style={-{Latex[length=2.2mm]}, thick, draw=borrador2},
      titulo/.style={font=\small\bfseries, text=borrador2, anchor=east},
      nota/.style={font=\footnotesize, text=borrador2, align=center},
    ]
      \node[titulo] at (-0.4,2.7) {Guías};
      \node[nodo] (guias) at (1.7,2.7) {Deduplicar por\\(nombre, contenido)};
      \node[nodo] (gproceso) at (5.8,2.7) {Añadir encabezado,\\trocear y reequilibrar};
      \node[resultado] (rguia) at (10.0,2.7) {guia};

      \node[titulo] at (-0.4,0) {Contenido directo};
      \node[nodo] (directo) at (1.7,0) {Salidas y asignaturas\\sin guía};
      \node[nodo] (dproceso) at (5.8,0) {Normalizar, encabezar\\y aplicar límites};
      \node[resultado] (rdirecto) at (10.0,0) {salidas\\asignatura\_sin\_guia};

      \node[titulo] at (-0.4,-2.7) {Vistas derivadas};
      \node[nodo] (derivadas) at (1.7,-2.7) {Construir resúmenes\\deterministas};
      \node[nodo] (vproceso) at (5.8,-2.7) {Añadir metadatos y\\aplicar límites};
      \node[resultado, minimum width=4.3cm] (rderivada) at (10.0,-2.7)
        {plan\_de\_estudios · catalogo\\ficha\_titulacion · mencion};

      \node[resultado, minimum width=4.0cm, minimum height=1.3cm] (json) at (14.7,0)
        {chunks.json\\{\scriptsize siete valores de origen}};

      \draw[flecha] (guias) -- (gproceso);
      \draw[flecha] (gproceso) -- (rguia);
      \draw[flecha] (directo) -- (dproceso);
      \draw[flecha] (dproceso) -- (rdirecto);
      \draw[flecha] (derivadas) -- (vproceso);
      \draw[flecha] (vproceso) -- (rderivada);
      \draw[flecha] (rguia) -- (json);
      \draw[flecha] (rdirecto) -- (json);
      \draw[flecha] (rderivada) -- (json);

      \node[nota, anchor=north] at (7.2,-4.0)
        {Máximo duro: 900 caracteres · mínimo preferido: 200 caracteres\\
         Ningún fragmento mezcla dos asignaturas; se identifican por (grado, código o nombre).};
    \end{tikzpicture}%
  }
  \caption[Alternativa del proceso de fragmentación]{\textcolor{borrador2}{Versión alternativa del proceso de fragmentación. Separa las guías, el contenido directo y las vistas derivadas; las tres vías convergen en una colección con siete orígenes y un máximo duro de 900 caracteres. Fuente: Elaboración propia.}}
  \label{fig:secuencia_chunking_alt}
\end{figure}
\endgroup
